\begin{table}[!htbp]
\caption{Definitions of the candidate HIs.}
\label{tab:2-2}
\centering
\zihao{6}
\setlength{\tabcolsep}{2.0pt}
\renewcommand{\arraystretch}{1.12}
\begin{ThreePartTable}
\begin{tabularx}{0.98\linewidth}{
>{\centering\arraybackslash}p{0.045\linewidth}
>{\centering\arraybackslash}p{0.105\linewidth}
>{\raggedright\arraybackslash}X
>{\centering\arraybackslash}p{0.045\linewidth}
>{\centering\arraybackslash}p{0.105\linewidth}
>{\raggedright\arraybackslash}X}
\toprule
No. & Category & Health indicator & No. & Category & Health indicator \\
\midrule
HI1 & CVT & MS-CCCT 标定电压区间内的恒流充电时间
& HI9 & DTV & 微分温度--电压曲线的峰谷差 \\
HI2 & IC & 增量容量曲线的峰值
& HI10 & DTC & 微分温度--容量曲线的峰值 \\
HI3 & IC & 增量容量峰值对应的电压
& HI11 & DTC & 微分温度--容量曲线峰值对应的容量 \\
HI4 & IC & 增量容量特征峰的半峰宽
& HI12 & DTC & 微分温度--容量特征峰的半峰宽 \\
HI5 & DTV & 微分温度--电压曲线的峰值
& HI13 & 窗口放电容量 & 电压范围 3.80--3.40 V 内的放电容量 \\
HI6 & DTV & 微分温度--电压曲线峰值对应的电压
& HI14 & 窗口放电容量 & 电压范围 3.20--3.00 V 内的放电容量 \\
HI7 & DTV & 微分温度--电压曲线的谷值
& HI15 & 能量效率 & 单次循环放电能量与充电能量之比 \\
HI8 & DTV & 微分温度--电压曲线谷值对应的电压
& & & \\
\bottomrule
\end{tabularx}
\end{ThreePartTable}
\end{table}
